\chapter{Вычислительный эксперимент}
\label{ch:experiments}

Метод проверяется на двух взаимодополняющих задачах: (i) контролируемой
двумерной модельной задаче, позволяющей проверить, совпадает ли
восстановленная маршрутизация с истинной кластерной структурой, и
(ii) разнородном бенчмарке на основе смеси CIFAR-100 и SVHN, где
проводится сравнение с базовыми методами.

\section{Модельная задача}
\label{sec:toy}

Сэмплируются $2000$ точек в $\mathbb{R}^2$ из задачи на четыре класса,
состоящей из одного «линейного» облака (два класса, разделённые
диагональю) и одного «кольцевого» облака (два класса --- внутреннее и
внешнее кольцо). Данные делятся $80/20$ на обучение и валидацию и
кластеризуются на $M = 20$ кластеров методом KMeans. Берётся $K = 2$
эксперта, так что идеальное решение однозначно: один эксперт должен
специализироваться на линейном облаке, другой --- на кольцевом.
Пространство поиска --- небольшое cell-пространство над двумерными
входами, чьё множество операций намеренно подобрано так, чтобы два
облака предпочитали \emph{разные} архитектуры. В частности, оно
содержит \emph{линейную} операцию
$\boldsymbol{x}\mapsto W\boldsymbol{x}$, которой уже достаточно для
разделения линейно-разделимого облака (два класса делятся прямой
диагональю), и операцию \emph{радиальной базисной функции} (RBF)
$\boldsymbol{x}\mapsto \exp\!\bigl(-\gamma\,\lVert\boldsymbol{x}-\boldsymbol{c}\rVert^2\bigr)$,
кодирующую расстояние до центра $\boldsymbol{c}$ и естественно подходящую
кольцевому облаку, чьи два класса (внутреннее и внешнее кольцо)
разделяются радиусом, а не гиперплоскостью; несколько вспомогательных
операций (тождество, поэлементный квадрат, сдвиг, $\mathrm{ReLU}$,
dropout) дополняют множество. Поскольку чисто линейная архитектура не
моделирует концентрические кольца, а RBF-доминированная избыточно плоха
на линейном облаке, оптимальная смесь вынуждена назначить
\emph{архитектурно различных} экспертов двум областям; тем самым
модельная задача проверяет именно архитектурную специализацию, а не
только маршрутизацию.

С параметрами, приведёнными в журнале экспериментов, SGEM восстанавливает
\emph{идеальное} разбиение кластеров $(10, 10)$ между двумя экспертами за
$20$ EM-итераций, совпадающее с эталонным разбиением по координате
центроида (рис.~\ref{fig:toy}). Эта проверка подтверждает, что
суррогатный объектив~\eqref{eq:moe_objective} корректно поставлен и что
предложенный оптимизатор решает его при достаточно точном суррогате, в
согласии с Теоремой~\ref{thm:sgem}.

\begin{figure}[h]
  \centering
  \begin{subfigure}[b]{0.32\textwidth}
      \centering
      \includegraphics[width=\textwidth]{toy_dataset}
      \caption{}
      \label{fig:toy_dataset}
  \end{subfigure}
  \hfill
  \begin{subfigure}[b]{0.32\textwidth}
      \centering
      \includegraphics[width=\textwidth]{toy_clusters}
      \caption{}
      \label{fig:toy_clusters}
  \end{subfigure}
  \hfill
  \begin{subfigure}[b]{0.32\textwidth}
      \centering
      \includegraphics[width=\textwidth]{toy_sgem_split}
      \caption{}
      \label{fig:toy_sgem}
  \end{subfigure}
  \caption{Модельная задача. (a) Датасет на четыре класса: линейное облако
  (классы 0--1) и кольцевое облако (классы 2--3). (b) $M = 20$
  KMeans-кластеров (кресты --- центроиды). (c) Разбиение экспертов,
  восстановленное SGEM при $K = 2$: кластеры делятся точно по границе
  линейное/кольцевое облако, совпадая с эталонным разбиением с точностью
  до перестановки меток экспертов.}
  \label{fig:toy}
\end{figure}

\section{Смесь CIFAR-100 и SVHN}
\label{sec:cifar}

Чтобы проверить, восстанавливает ли кластерно-зависимый объектив
осмысленное разбиение данных на реалистичных данных, строится разнородный
бенчмарк смешиванием CIFAR-100 и SVHN в равной пропорции (всего
$42\,000$ изображений). Метки SVHN сдвинуты на $100$, так что
объединённая задача содержит $110$ классов. Два источника образуют
максимально разделённую доменную структуру, задающую корректный эталон:
при $K = 2$ идеальная маршрутизация приписывает каждый кластер эксперту,
специализированному на его источнике. Данные разбиваются на $M = 30$
\emph{семантических} кластеров, полученных предобученным экстрактором
признаков, с трёхчастным разбиением обучение/валидация/тест
($0.7/0.15/0.15$). Итоговый эталон содержит $19$ CIFAR-мажоритарных и
$11$ SVHN-мажоритарных кластеров со средней чистотой источника $0.978$.

Все конфигурации MoE используют $K = 2$ эксперта. Базовая сеть --- малая
DARTS-подобная архитектура \texttt{CIFAR100Net} (stem, две обычные
ячейки, фиксированная редуцирующая ячейка, global average pooling,
полносвязная голова). MoE-обёртка инстанцирует $K$ таких сетей,
объединённых фиксированным жёстким \emph{кластерным шлюзом}: каждый
кластер направляется одному эксперту, так что шлюз напрямую отражает
обнаруженное разбиение. Поиск архитектуры использует разбиение
обучение/валидация; итоговая модель дообучается $100$ эпох на
$\text{обучение} \cup \text{валидация}$ и оценивается \emph{однократно}
на отложенном тесте.

\textbf{Балансировка нагрузки.} Без регуляризации динамика EM склонна
схлопывать все кластеры на одного эксперта --- это свойство самого
объектива. При жёсткой маршрутизации, когда каждый кластер $m$
направляется единственному эксперту $k(m)$, внутренняя сумма
в~\eqref{eq:moe_objective} имеет лишь одно слагаемое, и логарифм гасит
экспоненту:
\begin{equation}
\label{eq:hard_objective}
\mathcal{J} \;=\; -\sum_{k=1}^{K} N_k\, u_k,
\qquad N_k = \!\!\sum_{m:\,k(m)=k}\!\! |\mathcal{C}_m|,
\end{equation}
где $u_k = u(\boldsymbol{\alpha}_k, \mathbf{r}_k)$ --- предсказанная
ошибка эксперта $k$, а $N_k$ --- число направленных ему объектов. Если
суррогат почти постоянен по своему аргументу-маршрутизации и эксперты
взаимозаменяемы ($u_k \equiv \bar u$), то $\mathcal{J} = -\bar u\, N$ не
зависит от разбиения $\{N_k\}$; а если $u_k$ различны, то $\mathcal{J}$
линейна по $N_k$ и максимизируется отправкой всех кластеров одному
эксперту с наименьшим $u_k$. Чтобы этому противодействовать, M-шаг
дополняется штрафом балансировки
\begin{equation}
\label{eq:loadbalance}
\mathrm{LB}(\mathbf{R}) \;=\; K \cdot \sum_{k=1}^{K}\bar{r}_k \cdot p_k,
\qquad
\bar{r}_k = \frac{1}{M}\sum_{m=1}^{M} r_{mk}, \quad
p_k = \frac{|\{m : \arg\max_{k'} r_{mk'} = k\}|}{M},
\end{equation}
с весом $\lambda_{\text{LB}}$; он минимален при равномерной маршрутизации.
Поскольку правильное доменное разбиение само несбалансировано ($19/11$),
агрессивный штраф уводил бы маршрутизацию \emph{от} него, поэтому
$\lambda_{\text{LB}}$ удерживается малым ($\lambda_{\text{LB}} = 1$), а
схлопывание предотвращается схемой сбора данных, описанной ниже.
Приложение~\ref{app:lambda} делает этот компромисс количественным,
локализуя точку перелома $\lambda_{\text{LB}} \approx 3.77$, за которой
балансированное --- и потому доменно-слепое --- разбиение уже выигрывает
по пенализованному объективу.

\textbf{Честный сбор данных.} Ранняя версия пайплайна достигала лишь
$\approx\!0.59$ точности при почти случайном восстановлении источника
($\approx\!60\%$), что было вызвано двумя дефектами построения датасета
суррогата, а не оптимизатором EM. Во-первых, каждая
кандидат-архитектура оценивалась на \emph{всей} валидации, а не на
области своих кластеров, поэтому эксперт, обученный на чистом SVHN,
оценивался частично на изображениях CIFAR и выглядел плохим; это
схлопывало разрыв предсказанных ошибок между специализированным и
смешанным экспертом. Ограничение оценки кластерами самого эксперта
(оценка по области) устраняет смещение. Во-вторых, кластерные маски для
обучения суррогата сэмплировались включением каждого кластера независимо
с вероятностью $1/2$, так что почти чистые односторонние подмножества
практически не возникали; замена этого \emph{логравномерным} сэмплером по
размеру маски делает малые, часто почти чистые подмножества частыми ---
без обращения к меткам источника. С обоими исправлениями и
$\lambda_{\text{LB}} = 1$ SGEM восстанавливает доменную структуру
самостоятельно.

Результаты приведены в таблице~\ref{tab:cifar_main}. При честном сборе
данных SGEM достигает точности $0.6199 \pm 0.009$, превосходя как
DARTS$\times K$ ($0.5913$), так и базовый Random-MoE ($0.5467$) на
$+2.9$ и $+7.3$ процентных пункта, и оказываясь в пределах $1.3$ п.п. от
оракульного MoE ($0.6327$), которому \emph{дано} истинное разбиение по
источнику. Именно кластерный шлюз выделяет качество маршрутизации: при
обучаемом шлюзе методы становятся почти неразличимы (все $\approx\!0.63$),
так как обучаемый шлюз способен компенсировать неоптимальное разбиение.
Более тонкий ablation, раскладывающий эту точность на вклад
маршрутизации и архитектуры, приведён в Приложении~\ref{app:decomp}:
выигрыш почти полностью обусловлен маршрутизацией кластер$\to$эксперт,
тогда как архитектурная гетерогенность экспертов при уже правильной
маршрутизации добавляет мало, а обучение маршрутизации внутри SGEM (а не
её фиксация) как раз и восстанавливает доменную структуру.

\begin{table}[h]
\centering
\caption{Точность на тесте для смеси CIFAR-100 / SVHN с кластерным шлюзом
($M=30$, $K=2$, семантическая кластеризация, среднее $\pm$ стд.).
Оракульный MoE использует истинное разбиение по источнику и является
верхней границей, а не конкурентом.}
\label{tab:cifar_main}
\begin{tabular}{lcc}
\toprule
Метод & точность (кластерный шлюз) & сидов \\
\midrule
Random-MoE & $0.5467 \pm 0.029$ & 4 \\
DARTS$\times K$ (одна арх.) & $0.5913 \pm 0.009$ & 5 \\
\textbf{SGEM (наш)} & $\mathbf{0.6199 \pm 0.009}$ & 3 \\
\midrule
Оракульный MoE \emph{(метки источника)} & $0.6327 \pm 0.004$ & 3 \\
\bottomrule
\end{tabular}
\end{table}

Поскольку эталонное разбиение по источнику резкое (средняя чистота
$0.978$), можно напрямую проверить, согласуется ли с ним обученная
маршрутизация. По трём сидам жёсткое назначение SGEM теперь совпадает с
эталонным доменным разбиением на $92\% \pm 2\%$ кластеров (наилучшая из
двух перестановок меток) --- намного выше $\approx\!54\%$ сбалансированного
случайного разбиения и значительно лучше $\approx\!60\%$ ранней смещённой
версии. В типичном прогоне (рис.~\ref{fig:svhn_reps}) один эксперт
становится \emph{чистым} CIFAR-специалистом ($17$ кластеров, без SVHN), а
другой собирает все $11$ SVHN-кластеров вместе лишь с $2$ CIFAR. Таким
образом, кластерно-зависимый объектив \emph{действительно} обнаруживает
доменную структуру данных на этом бенчмарке без использования меток
источника.

\textbf{Сходимость.} На рис.~\ref{fig:convergence} прослежен один прогон
SGEM по EM-итерациям. Взвешенный по размерам объектив $L^{(t)}$ растёт с
точностью до шума оптимизации, а доля совпадения с источником
поднимается с $67\%$ до $100\%$, эмпирически иллюстрируя сходимость,
гарантированную Теоремой~\ref{thm:sgem}. Пики двух кривых расходятся на
одну итерацию (объектив на $t = 8$, восстановление на $t = 9$) --- тот же
рассинхрон объектива и качества у оптимума, что отмечен выше.

\begin{figure}[h]
    \centering
    \includegraphics[width=0.7\textwidth]{convergence_curve}
    \caption{Сходимость SGEM на смеси CIFAR-100/SVHN (seed $322$). Левая
    ось (синяя): совпадение жёсткой маршрутизации с эталонным разбиением
    по источнику; пунктир --- уровень случайного разбиения $54\%$. Правая
    ось (красная): взвешенный по размерам объектив $L^{(t)}$. Обе
    величины растут по EM-итерациям, доля совпадения достигает $100\%$ на
    $t = 9$.}
    \label{fig:convergence}
\end{figure}

\begin{figure}[h]
    \centering
    \includegraphics[height=0.42\textheight]{svhn_cluster_reps}
    \caption{Представители кластеров смеси CIFAR-100 / SVHN при
    маршрутизации SGEM ($K=2$; показаны первые $15$ из $M = 30$
    семантических кластеров). Каждая строка --- один кластер: левый
    столбец --- медоид кластера, остальные --- восемь ближайших к
    среднему изображений. Рамка медоида окрашена по эксперту, которому
    SGEM приписывает кластер (E0 --- SVHN-специалист, красный; E1 ---
    CIFAR-специалист, зелёный). SGEM направляет все SVHN-кластеры
    (изображения цифр) эксперту E0 и сохраняет E1 чистым
    CIFAR-экспертом; доля совпадения с источником для этого прогона
    --- $93\%$.}
    \label{fig:svhn_reps}
\end{figure}
